\section{What the certificates separate and refuse}
\label{sec:separations}

The catalog has 36 entries. This section states what they establish, sorted by kind:
separations exhibited in the reference world, statements that follow from the definitions,
refusals of credit, two refuted proposals and two deferred questions. The supplement
states each entry (\Cref{supp:register}).

\subsection{Separations}
\label{sec:sep-separations}

The main use of the language is to keep apart properties that are easily conflated. For 27
such pairs --- sound versus reachable, contact versus strict join, enablement versus descent,
route residue versus an arrow, and so on --- the finite reference world of \Cref{sec:ttw}
contains a record that has the first property and lacks the second. These separations hold
between certificate fields, so they show only that the fields are independent as records;
\Cref{tab:separations} lists them.

Four of them have semantic counterparts proved in this paper.
\Cref{prop:refinement} explains why an observable that factors globally through the pairing gains
nothing from a common refinement, and \Cref{rem:refinement-limit} states what strictness
leaves open.
\Cref{prop:strict-not-cheap} explains why relabelling produces no strict content.
The seed-partition row is a record that stipulates two seed partitions and stores unequal
terminal-package fields. It is a different statement from \Cref{ex:disabling}, which proves
that two admission orders from the same initial state reach different terminal states
because one admission disables another. When admission only enables,
\Cref{thm:admission} gives a unique terminal state from each state. The enablement row also
has a semantic counterpart (\Cref{prop:enablement-model}).

\subsection{Semantic models}
\label{sec:sep-semantic}

Four of the entries above are also proved about explicit transition systems rather than about
records.

\begin{proposition}[Enablement without descent]
\label{prop:enablement-model}
Let the lower theory have states $\{0,1,2\}$ and a step $a\to a'$ whenever $a<a'$. Form an upper
state from each lower state $a$ as the vector of states reachable from $a$, and let the upper
theory have a step between two formed states exactly when a lower step joins their seeds. Then
\begin{enumerate}
  \item the three formed states are distinct, and an upper step exists exactly when a lower step
        supplies it; in particular the formed state of $0$ steps to the formed state of $1$;
  \item the observable ``$0$ is reachable'' does not factor on its image through the quotient
        ``$1$ is reachable'': the formed states of $0$ and $1$ agree on the quotient and differ
        on the observable.
\end{enumerate}
An observable that is constant on the fibres of a quotient does factor through it.
\end{proposition}

\begin{proof}
Formation is injective because the vector of $a$ has first true entry at $a$; an upper step
between formed states is by definition a lower step between their seeds. The formed states of
$0$ and $1$ are $(1,1,1)$ and $(0,1,1)$, a split pair for the quotient and the observable, so
\Cref{thm:split-criterion} applies. The last sentence is the converse direction of the same
theorem.
\end{proof}

\Cref{fig:enablement} shows the two layers and the split pair.

\begin{figure}[htbp]
\centering
\begin{tikzpicture}[font=\footnotesize,>=Stealth,
  upper/.style={draw=gray,rounded corners=2pt,minimum width=2.55cm,minimum height=1.05cm,align=center,inner sep=3pt},
  lower/.style={draw=gray,rounded corners=2pt,minimum width=0.8cm,minimum height=0.6cm}]
  \node[anchor=east] at (0.65,0) {upper};
  \node[anchor=east] at (0.65,-3.3) {lower};
  \node[upper,very thick,draw=black] (u0) at (2.4,0)
    {$(1,1,1)$\\$q=1,\ f=1$};
  \node[upper,very thick,draw=black] (u1) at (6.5,0)
    {$(0,1,1)$\\$q=1,\ f=0$};
  \node[upper] (u2) at (10.6,0)
    {$(0,0,1)$\\$q=0,\ f=0$};
  \node[lower] (l0) at (2.4,-3.3) {$0$};
  \node[lower] (l1) at (6.5,-3.3) {$1$};
  \node[lower] (l2) at (10.6,-3.3) {$2$};
  \draw[->] (u0) -- (u1);
  \draw[->] (u1) -- (u2);
  \draw[->] (u0.north) to[out=65,in=115] (u2.north);
  \draw[->] (l0) -- (l1);
  \draw[->] (l1) -- (l2);
  \draw[->] (l0.south) to[out=-65,in=-115] (l2.south);
  \draw[->,dashed,gray] (l0) -- (u0);
  \draw[->,dashed,gray] (l1) -- (u1);
  \draw[->,dashed,gray] (l2) -- (u2);
\end{tikzpicture}
\caption{Lower steps and their induced upper steps between formed reachability vectors. Dashed arrows map each lower state to its vector; the outlined vectors agree on $q$ (``$1$ reachable'') and differ on $f$ (``$0$ reachable'').}
\label{fig:enablement}
\end{figure}


The model is small, but it separates the two properties as statements about transitions, not
as independent record fields: the upper step is enabled by the lower dynamics, and what the
upper theory observes is still not a function of the lower quotient.

\begin{proposition}[No free join]
\label{prop:no-free-join}
Consider events that require and produce capabilities, and call a finite sequence of events
lawful when every capability an event requires is initially available or produced by an earlier
event. If a join $j$ requires a capability $c$ that is not initially available, and $j$ is the
only event that produces $c$, then $j$ occurs in no lawful sequence. If every capability that $j$
requires is initially available, the sequence consisting of $j$ alone is lawful.
\end{proposition}

\begin{proof}
Induction on the lawful sequence. When $j$ is appended, $c$ must be initially available, which
it is not, or produced by an earlier event, which would be an earlier occurrence of $j$; the
induction hypothesis excludes that. The second statement is immediate from the definition of a
lawful sequence; the formal development checks it for a one-event system in which the join's
capability is initially supplied.
\end{proof}

\begin{proposition}[Selection creates no fact]
\label{prop:selection}
Let a state be a set of lower facts, and let a selection step replace a state $S$ by
$S \cap K$ for some set $K$. After any finite sequence of selection steps, every fact present
was present at the start. A step that inserts a fact into the empty state is not a selection
step.
\end{proposition}

\begin{proof}
$S \cap K \subseteq S$, and inclusion is transitive along the sequence.
\end{proof}

The proposition is immediate once selection is modelled as intersection; its use is to fix what
a pure downward selection is allowed to do.

\begin{proposition}[Bounded search]
\label{prop:search}
Let a finite list of states be scanned for a decidable target, and let the reachable set be the
states reachable from a start within $h$ steps.
\begin{enumerate}
  \item If every listed state is reachable within $h$ steps and the scan finds the target, some
        state reachable within $h$ steps satisfies it.
  \item If every state reachable within $h$ steps is listed and the scan finds nothing, no state
        reachable within $h$ steps satisfies the target; if moreover every reachable state is
        reachable within $h$ steps, no reachable state does.
\end{enumerate}
\end{proposition}

\begin{proof}
Both parts compare what the scan inspects with what is reachable: the first uses that every listed state is
reachable, the second that every reachable state is listed.
(1) The scan returns a listed state satisfying the target. (2) A reachable target state would
be listed, so the scan would find it.
\end{proof}

A negative search result is thus exactly as strong as the coverage of its list and the
closure of its horizon. For two Boolean components that each step from false to true, one step
at a time, the joint state $(\text{true},\text{true})$ is not reachable within one step and is
reachable within two; this is a rendezvous threshold, not yet a model of a channel carrying
contact between two theories.

\subsection{Statements that follow from the definitions}
\label{sec:sep-definitional}

The following catalog entries are definitions, or statements that follow in a line from
the definition of the certificate they concern. The first is stated as a proposition
because the later refusals use it.

\begin{proposition}[Bootstrap]
\label{prop:bootstrap}
In a closed admission regime with no admitted seed and no reachable generator, no first
extension is authorized. Each of the three routes to authorization (an admitted seed, a
reachable generator, open external provision) authorizes a first extension in some
regime.
\end{proposition}

\begin{proof}
Authorization is defined as the disjunction of the three routes, and open external
provision requires the family to be open. The escapes are exhibited by three regimes, one
for each route.
\end{proof}

The content of the proposition is the list of routes, which fixes exactly what a first
extension in a closed regime needs. The other entries of this kind --- neutral seeds, endogenous
enablement, downward selection, horizons, certified non-interaction, composition of enablement,
confluence at a declared scope, join status, order residue, the join-created residual and the
cross-family corollaries --- are stated in \Cref{supp:definitional}; the four with semantic
counterparts are \Cref{prop:enablement-model,prop:no-free-join,prop:selection,prop:search}.

\subsection{Refusals}
\label{sec:sep-refusals}

Eleven entries are refusals: a credit is withheld when a field its definition requires is
absent, and named routes restore it. \Cref{tab:refusals} lists them. In each case the
eligibility condition for the credit contains the missing field, so the refusal is
immediate from the definition. The content is the list of escape routes, each of which is
exhibited by a configuration that earns the credit. One refusal, occurrence from
reachability, is also accompanied by an explicit counterprofile.

{\small
\begin{longtable}{@{}>{\raggedright\arraybackslash}p{4.4cm}>{\raggedright\arraybackslash}p{5.2cm}>{\raggedright\arraybackslash}p{5.3cm}@{}}
\caption{Refusals of credit and their escape routes.}
\label{tab:refusals}\\
\toprule
\textbf{Refused credit} & \textbf{Missing field} & \textbf{Escape routes} \\
\midrule
\endfirsthead
\toprule
\textbf{Refused credit} & \textbf{Missing field} & \textbf{Escape routes} \\
\midrule
\endhead
\bottomrule
\endfoot
first extension in a closed regime & seed or reachable generator & admitted seed; reachable generator; open external provision \\
independence-sensitive join from resemblance & witnessed independent contact & witnessed independent contact \\
retrospective self-certification & prospective certificate & prospective certificate \\
total-lens transfer to a partial target & certified adapter & certified adapter \\
native credit with an unpriced observer & charged occupancy & external observer; zero occupancy \\
join from contact alone & composite, retention and anti-product witness & complete strict-join certificate \\
strictness from a product or common refinement & anti-product witness & anti-product witness \\
source independence from one lineage & disjoint roots & disjoint-root certificate \\
positive-cost join without payment & payment & paid channel (a certified zero-cost channel credits only a join of amount zero) \\
arrow from route residue & drive certificate & independently driven arrow \\
occurrence from reachability & firing & a fired occurrence \\
\end{longtable}
}

Three further refusals are worth recording with them. A join cannot create the first
capability needed to execute itself, cannot manufacture strict content by relabelling, and
cannot receive positive-cost formation for nothing. The second follows from
\Cref{prop:strict-not-cheap} and the third is the payment refusal above. The first is
\Cref{prop:no-free-join}: a capability that only the join produces cannot authorize the
join's own execution.

An observer or instrument that uses bounded shared capacity cannot certify native or
endogenous formation while its occupancy is under-charged or its source role or cost is
off the ledger. Observation is not refused; unrecorded observation is.

\subsection{Two refuted proposals}

\begin{remark}[Refinement monotonicity is false]
\label{rem:refinement-refuted}
The unconditional claim that refining a parent preserves the join is refuted at the
level of records. A refinement record carries an effect label (preserve, strengthen, weaken
or destroy) and preservation flags. A well-formed record labelled \emph{destroys} exists,
so the certificate language does not permit the unconditional claim. What survives is a
classification of refinement records by their label, and all four labels have well-formed record
witnesses. As a contrast in observable factorization, not a model of refinement acting on a join,
the three-bit join shows both behaviours: the pairing, a coarse common refinement, loses the
observable $h$, while the identity refinement retains it (\Cref{rem:refinement-limit}). No model of refinement acting on a join is given.
\end{remark}

\begin{remark}[Two contact quantities disagree]
\label{rem:degree-refuted}
One might expect contact to carry a conserved scalar degree. A contact quantity records five
coordinates: access, retention, novelty, observer occupancy and residual. For the vector
$(1,1,1,0,1)$ the plain sum is $4$, while the sum that weights novelty by $2$ and residual by
$3$ is $7$, so the two plausible degrees disagree. An apparent conservation that
omits the observer's occupancy fails the accounting certificate, while an exact scoped
conservation certificate, with every coordinate included, is available. No universal
statement about conserved quantities is proved in either direction: the two scalars
disagree, and the accounting stays local.
\end{remark}

\subsection{Two deferred questions}

\begin{remark}[Categorical reduction]
\label{rem:deferral-categorical}
Whether a join of theories reduces to a product, pullback or pushout
\citep{MacLane1998} is not settled here. The catalog's categorical certificate is a
checklist of flags (identities checked, composition closed, associativity checked, a
universal property recorded) with no morphisms or composition, so it supports no
categorical statement. What is proved is the set-level fact of \Cref{thm:strict-split}: a
strict observable does not factor, on its image, through the pairing $J\to A\times B$, which is the map
into the product of sets. A categorical treatment would need hom-sets, composition and a
universal property quantified over mediating maps.
\end{remark}

\begin{remark}[Primitive-operation algebra]
\label{rem:deferral-algebra}
A generators-and-relations presentation of the primitive operations is not attempted. The
part the paper uses is the monoid of resource deltas described above.
\end{remark}
